% ID: FIS-URT-ANE-001
% Titolo: Urto anelastico tra scatoloni collegati a una molla
% Disciplina: Fisica
% Area: Quantita di moto e urti
% Argomento: Urti anelastici
% Sottoargomento: Urto completamente anelastico con energia elastica
% Classe: Terza liceo scientifico
% Difficolta: 4
% Tipo: problema_numerico
% Risultato: x_max = 6,45 cm
% Tempo_stimato: 12 min
% Competenze: conservazione dell'energia meccanica, conservazione della quantita di moto, distinzione tra fasi del moto
% Prerequisiti: energia potenziale elastica, quantita di moto, urto completamente anelastico
% Tag: fisica, quantita_moto, urti_anelastici, molla, conservazione_energia
% Fonte: rielaborazione_da_traccia_fornita_dal_docente
% Autore: Riccardo Carli
% Licenza: CC-BY-SA-4.0

\begin{esercizio}
Due scatoloni di masse \(m_A=0{,}15\,\mathrm{kg}\) e
\(m_B=0{,}37\,\mathrm{kg}\) si trovano su un piano orizzontale privo di
attrito. Lo scatolone \(A\) e collegato a un muro mediante una molla ideale.
Quando la molla e alla lunghezza di equilibrio, \(A\) e a contatto con \(B\).

Si comprime la molla di \(s=12\,\mathrm{cm}\), spostando \(A\) verso il muro,
mentre \(B\) rimane fermo. Rilasciato \(A\), esso torna verso \(B\) e lo urta
proprio quando la molla raggiunge la lunghezza di equilibrio. L'urto e
completamente anelastico.

Determina il massimo spostamento verso destra del sistema formato dai due
scatoloni dopo l'urto.
\end{esercizio}

\begin{soluzione}
Indichiamo con \(k\) la costante elastica della molla. Nella prima fase, dal
rilascio fino all'urto, si conserva l'energia meccanica:
\[
\frac12ks^2=\frac12m_Av_A^2.
\]
Quindi
\[
v_A^2=\frac{k}{m_A}s^2.
\]

Durante l'urto, invece, l'energia cinetica non si conserva, perche l'urto e
completamente anelastico. Si conserva la quantita di moto:
\[
m_Av_A=(m_A+m_B)V,
\]
dove \(V\) e la velocita comune subito dopo l'urto. Dunque
\[
V=\frac{m_A}{m_A+m_B}v_A.
\]

Dopo l'urto i due scatoloni si muovono insieme e la molla si allunga. Fino al
punto di massimo allungamento non vi sono forze dissipative, quindi l'energia
meccanica si conserva nuovamente:
\[
\frac12(m_A+m_B)V^2=\frac12kx_{\max}^2.
\]
Sostituendo l'espressione di \(V\):
\[
x_{\max}^2
=\frac{m_A+m_B}{k}
\left(\frac{m_A}{m_A+m_B}\right)^2v_A^2.
\]
Usando
\[
v_A^2=\frac{k}{m_A}s^2
\]
si ottiene
\[
x_{\max}^2=\frac{m_A}{m_A+m_B}s^2,
\]
e quindi
\[
x_{\max}=s\sqrt{\frac{m_A}{m_A+m_B}}.
\]
Numericamente:
\[
x_{\max}=0{,}12
\sqrt{\frac{0{,}15}{0{,}15+0{,}37}}
\simeq0{,}0645\,\mathrm{m}.
\]
Pertanto
\[
\boxed{x_{\max}\simeq6{,}45\,\mathrm{cm}}.
\]
Lo spostamento e misurato dalla posizione dell'urto, che coincide con la
posizione di equilibrio della molla.
\end{soluzione}
